% =====================================================================
%  CHƯƠNG 6 — KẾT LUẬN VÀ HƯỚNG PHÁT TRIỂN
%  Nguồn facts: architecture/* + DECISIONS.md. Không bịa số đo.
% =====================================================================
\documentclass[../DoAn.tex]{subfiles}
\begin{document}

\section{Kết luận}
\label{section:6.1}

Đồ án đã xây dựng thành công một hệ thống mô phỏng mạng truyền thông trên ô tô
hoàn chỉnh theo mô hình ba tầng: Nút cảm biến (STM32F103) thu thập đầu vào và sở
hữu trạng thái xe, Nút cổng kết nối (STM32F103 + W5500) dịch giao thức từ CAN
sang Ethernet/TCP, và ứng dụng Android trên Raspberry Pi~4 hiển thị cụm đồng hồ
số. Toàn tuyến dữ liệu \textit{đầu vào vật lý $\to$ CAN $\to$ Gateway $\to$ TCP
$\to$ giao diện} đã vận hành thông suốt và được kiểm thử ở cả cấp module lẫn cấp
tích hợp (Chương~4).

\subsection*{Những kết quả đã đạt được}
\begin{itemize}
  \item Mạng CAN hai khung (\texttt{0x100} điều khiển, \texttt{0x200} cảm biến)
    ở 250~kbps, có kiểm lỗi CRC8 và theo dõi số thứ tự ở lớp ứng dụng.
  \item Cầu nối CAN--Ethernet thời gian thực với cơ chế đệm hàng đợi và socket
    không chặn, không mất khung khi phía nhận chưa sẵn sàng và tự phục hồi sau sự
    cố mạng (\ref{section:5.1}).
  \item Giao thức tải tin tối giản 18~byte tận dụng bảo đảm của TCP, kèm cơ chế
    tích luỹ trạng thái lai cho ảnh chụp luôn đầy đủ (\ref{section:5.2}).
  \item Ứng dụng hiển thị bền vững nhờ dịch vụ nền, tự khởi động lại và tự kết nối
    lại, với giao diện cụm đồng hồ tự co giãn theo độ phân giải màn hình.
\end{itemize}

\subsection*{So sánh với các sản phẩm tương tự}
So với một cụm đồng hồ số thương mại trên xe hiện đại, sản phẩm của đồ án là một
\emph{mô hình mô phỏng phục vụ học tập}, không nhằm thay thế hệ thống đạt chuẩn ô
tô. Các hệ thống thương mại thường vận hành trên nền tảng phần mềm chuẩn hoá
(ví dụ AUTOSAR), cổng kết nối đạt chuẩn an toàn chức năng và sử dụng CAN-FD hoặc
Automotive Ethernet. Tuy nhiên, trong phạm vi học thuật, đồ án tái hiện trung
thực \emph{triết lý kiến trúc} của một mạng ô tô thực: phân tách rõ ECU nguồn
trạng thái -- gateway dịch giao thức -- head-unit hiển thị, với ranh giới trách
nhiệm được giữ vững xuyên suốt. Đây là điểm tương đồng cốt lõi về mặt thiết kế,
đồng thời là giá trị giáo dục của sản phẩm.

\subsection*{Những phần chưa hoàn thiện}
\begin{itemize}
  \item Một số chỉ số hiệu năng (độ trễ đầu cuối, thông lượng tối đa, tỉ lệ mất
    gói, thời lượng soak test) chưa được đo đạc đầy đủ và còn để trống trong các
    kịch bản kiểm thử (Chương~4).
  \item Hệ thống dùng địa chỉ IP tĩnh và liên kết điểm-điểm, chưa hỗ trợ nhiều
    màn hình/khách hàng hay cấu hình mạng linh hoạt.
  \item Chưa có lớp bảo mật cho kênh truyền và chưa hỗ trợ các nghi thức chẩn đoán
    chuẩn ô tô.
\end{itemize}

\subsection*{Bài học kinh nghiệm}
Quá trình thực hiện đem lại nhiều bài học thực tiễn: (1) trong hệ nhúng đa nhiệm,
\emph{thứ tự khởi tạo} và \emph{quyền sở hữu tài nguyên} quyết định độ ổn định
nhiều như bản thân thuật toán (\ref{section:5.4}); (2) khi tầng dưới đã cung cấp
một bảo đảm, lặp lại nó ở tầng trên chỉ làm tăng độ phức tạp mà không tăng độ tin
cậy (\ref{section:5.2}); (3) việc giữ vững ranh giới trách nhiệm giữa các tầng
giúp khắc phục lỗi đúng chỗ và tránh lan toả thay đổi (\ref{section:5.6}); và
(4) log chẩn đoán có chủ đích là công cụ gỡ lỗi hiệu quả hơn nhiều so với thông
báo lỗi chung chung (\ref{section:5.5}).

\section{Hướng phát triển}
\label{section:6.2}

\subsection*{Hoàn thiện các chức năng hiện có}
\begin{itemize}
  \item \textbf{Đo đạc và đánh giá hiệu năng:} hoàn tất các phép đo còn để trống
    — độ trễ đầu cuối, thông lượng, tỉ lệ mất gói, độ ổn định khi chạy dài — kèm
    ảnh chụp dạng sóng CAN, ảnh giải mã và ảnh Wireshark để minh chứng.
  \item \textbf{Mở rộng tập tín hiệu:} bổ sung thêm khung/tín hiệu CAN (nhiệt độ,
    áp suất, quãng đường thực) và đưa các cảnh báo hiện đang suy diễn ở ứng dụng
    về đúng nguồn dữ liệu nếu cần độ chân thực cao hơn.
  \item \textbf{Cấu hình mạng linh hoạt:} hỗ trợ cấu hình địa chỉ động hoặc khám
    phá thiết bị, thay cho IP tĩnh cố định.
\end{itemize}

\subsection*{Các hướng nâng cấp mới}
\begin{itemize}
  \item \textbf{Nâng cấp lớp vật lý:} chuyển sang CAN-FD để tăng dung lượng tải
    và tốc độ, hoặc Automotive Ethernet cho băng thông cao hơn.
  \item \textbf{Chẩn đoán theo chuẩn ô tô:} tích hợp nghi thức chẩn đoán (ví dụ
    UDS) để đọc lỗi và giám sát tình trạng các nút.
  \item \textbf{Bảo mật truyền thông:} bổ sung xác thực và mã hoá kênh TCP, cũng
    như cơ chế chống giả mạo khung trên CAN.
  \item \textbf{Mở rộng quy mô hiển thị:} cho phép nhiều màn hình/khách hàng cùng
    nhận trạng thái, tiến tới mô hình phân phối dữ liệu nhiều người tiêu thụ.
  \item \textbf{Giám sát từ xa và cập nhật phần mềm:} đẩy dữ liệu telemetry lên
    nền tảng đám mây và hỗ trợ cập nhật firmware qua mạng (OTA).
\end{itemize}

\end{document}
